\scanpage{45}
le groupe distingué \deleted{qu'il} engendré par $\sigma$ \inserted{contient les rotations paires de $Z$} \deleted{donc est \unsure{indice}} \inserted{égal à $G$} si $n$ impair, et sinon \unsure{un|il} \deleted{\unsure{groupe}} \deleted{\unsure{se}} est isomorphe à $\{\pm1\}\mathbin{.}\mathbb Z/n'\mathbb Z$ (produit $\frac12$ direct) où $n'=n/2$ -- donc il est non cyclique, donc $\sigma$ n'est pas contenu dans un ss-groupe invariant cyclique $\neq G$, donc tous ces groupes sont contenus dans $Z$, qui est \inserted{donc} le plus grand de ces groupes.] Dans le cas de $\mathrm{Aff}(S)$\ldots.
% slanted
\margin{faire dém\ldots}

Donc sauf erreur, la structure quasi-polygonale combinatoire de $S\inserted{\subset E}$ est \inserted{uniquemt} déterminée, pour son système régulier de $n$ sommets dans $E$.

Principe de classification géométrique, pour $n$ fixé : Un système régulier de $n$ sommets (\deleted{\unsure{un|en}}\inserted{un} $\mathrm{SR}\deleted{\unsure{S}}n\mathrm S$) détermine

a) Un sous groupe cyclique d'ordre $n$
\[
Z\subset\mathrm{Aff}(E)
\]

b) Une orbite $\displaystyle S\in {}_Z\!\backslash E$

avec les conditions suivantes

1) $S$ non contenu dans une droite (cela implique que $Z$ opère \uline{fidèlemt} sur $S$, qui est donc un $Z$-torseur, $Z$ étant cyclique)

2) $\exists\,\sigma\in\operatorname{Norm}(Z)$ \deleted{$\unsure{\mathrm{Aff}|{}\neq\mathrm{id}_E}$} tel que $\sigma$ \deleted{\unsure{laisse fixe un pt de $S$|laisse invariant un pt de $S$}} \inserted{opère sur $Z$ par $z\mapsto z^{-1}$, et que $\sigma$ \unsure{invarie|inverse|invariant}} un pt de $S$ --
% archive sheet number: 44
